% 方板自由振动数值算例
% 格式对齐第四章：三线表（无竖线）、caption在表上方、不用\scriptsize、图片\linewidth

\section{数值算例}

通过数值算例系统检验不同边界条件下，混合格式在自由振动分析中的计算精度。本文数值算例采用无量纲材料参数，弹性模量 $E=1.0$，泊松比 $\nu=0.3$，密度 $\rho=1.0$，板厚 $h=10^{-3}$，板边长 $a=b=1.0$，剪切修正系数 $\kappa=5/6$。

无量纲频率参数定义为
\begin{equation}
\lambda = \frac{\bar{\omega}}{\pi^2}, \qquad \bar{\omega} = \omega a^2 \sqrt{\frac{\rho h}{D^b}}, \qquad D^b = \frac{E h^3}{12(1-\nu^2)}.
\end{equation}

参考解取自 Liew 等人\cite{liew1993} 采用 $pb$-$2$ Rayleigh-Ritz 法计算的结果，对应 $a/b=1.0$、$t/b=0.001$。

\subsection{频率对比}

五种数值方法用于对比：标准有限元（FEM）、缩减积分有限元（FEM缩减积分）、三变量混合有限元（FEM三变量）、无网格法（MF）、改进无网格法（MF2）。网格密度从 $20\times20$ 到 $30\times30$。

% TODO[数据口径]：当前FEM列为h=1e-2结果，待重跑为h=1e-3后更新
% TODO[方法说明]：MF与MF2最大相对差仅2.7e-4，非独立方法，需在表注说明
% TODO[对称性]：CSCS与SCSC在方板自由振动下由90°旋转对称严格等价，需合并或注明

\begin{table}[H]
\centering
\caption{CCCC边界前十阶无量纲频率 $\lambda$ 对比（Tri3网格，$25\times25$，$t/b=0.001$）}
\label{tab:vib_cccc}
\begin{tabular}{crrrrrrr}
\hline
阶次 & 标准解\cite{liew1993} & FEM & FEM缩减积分 & FEM三变量 & MF & MF2 \\
\hline
1 & 3.6463 & 9.0161 & 81.6867 & 3.6755 & 3.6466 & 3.6466 \\
2 & 7.4364 & 16.7062 & 148.2861 & 7.5235 & 7.4383 & 7.4385 \\
3 & 7.4364 & 19.7881 & 181.7857 & 7.5248 & 7.4383 & 7.4386 \\
4 & 10.9647 & 25.6354 & 226.2437 & 11.1483 & 10.9676 & 10.9682 \\
5 & 13.3315 & 32.0603 & 288.7722 & 13.5552 & 13.3378 & 13.3378 \\
6 & 13.3951 & 34.7997 & 321.2210 & 13.6267 & 13.4009 & 13.4009 \\
7 & 16.7177 & 36.6469 & 321.2819 & 17.0684 & 16.7241 & 16.7245 \\
8 & 16.7177 & 44.0710 & 396.6453 & 17.0803 & 16.7242 & 16.7248 \\
9 & 21.3296 & 49.6168 & 433.3729 & 21.8123 & 21.3440 & 21.3437 \\
10 & 21.3296 & 52.0191 & 471.5575 & 21.8188 & 21.3441 & 21.3438 \\
\hline
\end{tabular}
\end{table}

\begin{table}[H]
\centering
\caption{SSSS边界前十阶无量纲频率 $\lambda$ 对比（Tri3网格，$25\times25$，$t/b=0.001$）}
\label{tab:vib_ssss}
\begin{tabular}{crrrrrrr}
\hline
阶次 & 标准解 & FEM & FEM缩减积分 & FEM三变量 & MF & MF2 \\
\hline
1 & 2.0000 & 5.2493 & 47.5683 & 2.0078 & 2.0000 & 2.0000 \\
2 & 5.0000 & 11.2975 & 99.6490 & 5.0177 & 5.0006 & 5.0008 \\
3 & 5.0000 & 13.9179 & 128.0485 & 5.0292 & 5.0006 & 5.0008 \\
4 & 8.0000 & 18.9447 & 165.9420 & 8.0589 & 8.0002 & 8.0009 \\
5 & 10.0000 & 24.3030 & 218.3133 & 10.0673 & 10.0040 & 10.0042 \\
6 & 10.0000 & 26.7288 & 247.1585 & 10.0778 & 10.0040 & 10.0042 \\
7 & 13.0000 & 28.5535 & 248.4308 & 13.1272 & 13.0019 & 13.0029 \\
8 & 13.0000 & 35.2044 & 315.4769 & 13.1377 & 13.0021 & 13.0031 \\
9 & 17.0000 & 40.1562 & 348.1101 & 17.1846 & 17.0135 & 17.0139 \\
10 & 17.0000 & 41.8796 & 379.1305 & 17.1972 & 17.0136 & 17.0141 \\
\hline
\end{tabular}
\end{table}

\begin{table}[H]
\centering
\caption{CSCS边界前十阶无量纲频率 $\lambda$ 对比（Tri3网格，$25\times25$，$t/b=0.001$）}
\label{tab:vib_cscs}
\begin{tabular}{crrrrrrr}
\hline
阶次 & 标准解\cite{liew1993} & FEM & FEM缩减积分 & FEM三变量 & MF & MF2 \\
\hline
1 & 2.9334 & 7.3545 & 66.6355 & 2.9543 & 2.9336 & 2.9336 \\
2 & 5.5465 & 13.4000 & 119.3601 & 5.5886 & 5.5476 & 5.5478 \\
3 & 7.0242 & 17.7497 & 162.2562 & 7.1026 & 7.0257 & 7.0260 \\
4 & 9.5833 & 22.0418 & 193.7378 & 9.7214 & 9.5848 & 9.5855 \\
5 & 10.3564 & 26.9177 & 244.5636 & 10.4559 & 10.3613 & 10.3614 \\
6 & 13.0797 & 32.2880 & 282.2141 & 13.2967 & 13.0851 & 13.0852 \\
7 & 14.2053 & 32.6540 & 299.0921 & 14.4385 & 14.2097 & 14.2105 \\
8 & 15.6815 & 39.7010 & 355.5860 & 15.9922 & 15.6854 & 15.6862 \\
9 & 17.2597 & 44.5333 & 388.0839 & 17.4838 & 17.2743 & 17.2746 \\
10 & 20.2442 & 44.5796 & 408.3700 & 20.7280 & 20.2487 & 20.2499 \\
\hline
\end{tabular}
\end{table}

% TODO[对称性]：SCSC与CSCS在方板自由振动下由90°旋转对称严格等价，数据完全相同，可考虑合并

\begin{table}[H]
\centering
\caption{FSCS边界前十阶无量纲频率 $\lambda$ 对比（Tri3网格，$25\times25$，$t/b=0.001$）}
\label{tab:vib_fscs}
\begin{tabular}{crrrrrrr}
\hline
阶次 & 标准解\cite{liew1993} & FEM & FEM缩减积分 & FEM三变量 & MF & MF2 \\
\hline
1 & 1.2853 & 3.4320 & 31.1699 & 1.2888 & 1.2857 & 1.2860 \\
2 & 3.3502 & 8.2615 & 72.5739 & 3.3689 & 3.3492 & 3.3495 \\
3 & 4.2253 & 11.1267 & 102.6841 & 4.2428 & 4.2276 & 4.2293 \\
4 & 6.3846 & 15.0915 & 131.8066 & 6.4329 & 6.3852 & 6.3867 \\
5 & 7.3352 & 19.4851 & 176.4565 & 7.3935 & 7.3344 & 7.3349 \\
6 & 9.1817 & 22.6276 & 201.0210 & 9.2476 & 9.1893 & 9.1925 \\
7 & 10.4521 & 24.2808 & 218.1133 & 10.5718 & 10.4502 & 10.4519 \\
8 & 11.3372 & 29.9415 & 268.3376 & 11.4564 & 11.3437 & 11.3466 \\
9 & 13.3161 & 34.0026 & 296.0380 & 13.4776 & 13.3175 & 13.3180 \\
10 & 15.4787 & 35.0810 & 319.2567 & 15.7094 & 15.4811 & 15.4843 \\
\hline
\end{tabular}
\end{table}

\begin{table}[H]
\centering
\caption{FSSS边界前十阶无量纲频率 $\lambda$ 对比（Tri3网格，$25\times25$，$t/b=0.001$）}
\label{tab:vib_fsss}
\begin{tabular}{crrrrrrr}
\hline
阶次 & 标准解\cite{liew1993} & FEM & FEM缩减积分 & FEM三变量 & MF & MF2 \\
\hline
1 & 1.1847 & 3.1356 & 28.4668 & 1.1865 & 1.1841 & 1.1844 \\
2 & 2.8127 & 7.2761 & 64.2912 & 2.8232 & 2.8116 & 2.8119 \\
3 & 4.1743 & 10.6714 & 97.7653 & 4.1899 & 4.1763 & 4.1778 \\
4 & 5.9846 & 13.7398 & 119.9629 & 6.0196 & 5.9853 & 5.9868 \\
5 & 6.2679 & 17.6729 & 161.0025 & 6.2981 & 6.2669 & 6.2673 \\
6 & 9.1486 & 21.6411 & 188.6070 & 9.2128 & 9.1569 & 9.1600 \\
7 & 9.5730 & 23.2016 & 210.6866 & 9.6495 & 9.5716 & 9.5733 \\
8 & 11.0356 & 28.0091 & 250.4310 & 11.1336 & 11.0418 & 11.0447 \\
9 & 11.7212 & 31.6099 & 277.9959 & 11.8096 & 11.7231 & 11.7237 \\
10 & 14.7558 & 32.8124 & 296.9094 & 14.9263 & 14.7586 & 14.7617 \\
\hline
\end{tabular}
\end{table}

\subsection{收敛分析}

收敛分析采用绝对差值
\begin{equation}
e = |\lambda_{\mathrm{num}} - \lambda_{\mathrm{ref}}|,
\end{equation}
其中 $\lambda_{\mathrm{num}}$ 为数值解，$\lambda_{\mathrm{ref}}$ 为标准解。图\ref{vib_convergence} 给出六种边界条件下第一阶模态的收敛曲线。图中横坐标为 $\log_{10}(h)$，$h=1/n$ 为单元尺寸；纵坐标为 $\log_{10}(e)$，误差值越往下表示结果越精确。

\begin{figure}[H]
\centering
\includegraphics[width=0.9\textwidth]{vibration/convergence_first_mode_all.png}
\caption{六种边界条件第一阶模态无量纲频率的对数误差收敛曲线}
\label{vib_convergence}
\end{figure}

% TODO[图件中文化+去边框]：当前收敛图标签为英文且有坐标轴边框，需重画为无框图

\subsection{模态对比}

四边简支（SSSS）方板的自由振动模态具有解析形式，可直接作为振型对比的基准。图\ref{vib_modal_ssss} 给出 SSSS 边界下前六阶模态的变形图对比。每张图的列为不同数值方法，行为模态阶次。

% TODO[图件中文化+去边框]：当前模态图标签为英文且有坐标轴边框，需重画

\begin{figure}[H]
\centering
\textbf{(a) 第 1 阶}\\[1mm]
\includegraphics[width=\linewidth]{vibration/modal_shapes_SSSS.png}\\[3mm]
\caption{四边简支（SSSS）方板前六阶模态变形图对比}
\label{vib_modal_ssss}
\end{figure}

% TODO[待补充]：CCCC等其他边界的模态图（当前无标准模态，待后续补充）
% TODO[待补充]：收敛率定量值（MF、FEM三变量混合元对各边界的网格收敛斜率）
% TODO[待补充]：结果分析文字（读图三件套：图例解读→读法说明→结论）
